\documentclass[11pt]{article}
\usepackage[margin=1in]{geometry}
\usepackage{amsmath,amssymb,amsthm}
\usepackage{graphicx}
\usepackage[hidelinks]{hyperref}

\newtheorem{theorem}{Theorem}
\newtheorem{lemma}{Lemma}

\title{A Two-Atom Counterexample to the Nonlinear-Phase Exponential Basis Question}
\author{fa\_banach\_001 lane 1}
\date{2026-07-03}

\begin{document}
\maketitle

\begin{abstract}
Gabardo, Lai, and Oussa ask whether every finite Borel measure $\mu$ can
admit an orthogonal basis of generalized exponentials
$E(\Lambda,\varphi)=\{e^{2\pi i\lambda\cdot\varphi(x)}:\lambda\in\Lambda\}$
for some Borel phase $\varphi$. We give a literal counterexample: an unequal
two-atom probability measure. The obstruction is that generalized exponentials
have modulus one at every atom, so two such vectors can be orthogonal in a
two-atom weighted $L^2$ space only when the two atom weights are equal.
\end{abstract}

\section{Source Question}

The source paper is:

\begin{quote}
Jean-Pierre Gabardo, Chun-Kit Lai, and Vignon Oussa,
\emph{On exponential bases and frames with non-linear phase functions and some
applications}, arXiv:2007.03736.
\end{quote}

Section 8, ``Open problems'', asks:

\begin{quote}
Given any finite Borel measure $\mu$ on $\mathbb{R}^d$. Due to the flexibility
of Borel measurable functions $\varphi$, is it true that every $L^2(\mu)$ can
admit some $E(\Lambda,\varphi)$ as an orthogonal basis? Can we find a Borel
measure that does not admit any orthogonal basis of non-linear phase?
\end{quote}

The rendered source page is included in Figure~\ref{fig:source}.

\begin{figure}[h]
\centering
\includegraphics[width=0.78\textwidth]{figures/open_problem_crop.png}
\caption{Source open-problem page from arXiv:2007.03736, Section 8.}
\label{fig:source}
\end{figure}

\section{Definitions}

For a countable $\Lambda\subset\mathbb{R}^d$ and a Borel measurable phase
$\varphi$, the generalized exponential system is
\[
E(\Lambda,\varphi)=
\{x\mapsto e^{2\pi i\lambda\cdot\varphi(x)}:\lambda\in\Lambda\}.
\]
In an atomic $L^2(\mu)$ space, these are simply vectors with all coordinates
of modulus one.

\section{Proof Intuition}

The phase has enormous freedom, but it cannot change the weights of atoms.
For a measure with two atoms of unequal masses, $L^2(\mu)$ is a two-dimensional
weighted Hilbert space. An orthogonal basis would require two nonzero
orthogonal generalized exponentials. Since both vectors have coordinate
moduli equal to one, their inner product is a sum of two complex numbers whose
moduli are exactly the two atom masses. Two complex numbers can sum to zero
only if their moduli are equal. Unequal atom weights therefore prevent even a
pair of orthogonal generalized exponentials.

\begin{lemma}
Let $\mu=a\delta_{x_0}+b\delta_{x_1}$ with $a,b>0$ and $a\neq b$. For every
Borel measurable $\varphi:\{x_0,x_1\}\to\mathbb{R}^d$ and every countable
$\Lambda\subset\mathbb{R}^d$, the system $E(\Lambda,\varphi)$ contains no two
orthogonal nonzero vectors in $L^2(\mu)$.
\end{lemma}

\begin{proof}
Take two frequencies $\lambda,\lambda'\in\Lambda$. Their inner product is
\[
\left\langle e^{2\pi i\lambda\cdot\varphi},
e^{2\pi i\lambda'\cdot\varphi}\right\rangle_{L^2(\mu)}
=
a e^{2\pi i(\lambda-\lambda')\cdot\varphi(x_0)}
+b e^{2\pi i(\lambda-\lambda')\cdot\varphi(x_1)}.
\]
The two summands have moduli $a$ and $b$. If the inner product were zero, then
\[
a e^{2\pi i(\lambda-\lambda')\cdot\varphi(x_0)}
=
-b e^{2\pi i(\lambda-\lambda')\cdot\varphi(x_1)}.
\]
Taking absolute values gives $a=b$, contradicting the hypothesis. Thus no two
distinct generalized exponentials are orthogonal.
\end{proof}

\begin{theorem}
There exists a finite Borel measure $\mu$ on $\mathbb{R}$ such that no Borel
phase $\varphi$ and no countable frequency set $\Lambda$ make
$E(\Lambda,\varphi)$ an orthogonal basis for $L^2(\mu)$.
\end{theorem}

\begin{proof}
Let
\[
\mu=\frac13\delta_0+\frac23\delta_1.
\]
Then $L^2(\mu)$ is two-dimensional. Any orthogonal basis of this Hilbert space
must contain two nonzero orthogonal vectors. By the lemma, no generalized
exponential system $E(\Lambda,\varphi)$ contains such a pair. Hence no such
system is an orthogonal basis for $L^2(\mu)$.
\end{proof}

\section{Verification Notes}

The proof is independent of regularity of the phase; it applies to every Borel
phase and every countable frequency set. It also covers linear phases, so it
is stronger than merely excluding nonlinear phases.

Equivalently, using the source's Theorem 3.2, an essential-injective phase
would push the two unequal atom weights to a two-atom measure with the same
unequal weights. A finite atomic spectral measure with two atoms must have
equal weights, since an orthogonal pair of modulus-one vectors in the
two-dimensional weighted space exists only in that case.

\section{Novelty Check}

Before promotion, the run searched the local registry for arXiv:2007.03736 and
the title/key phrases, finding no existing packet for this question. Web
queries used the exact open-question wording and close variants involving
``orthogonal basis of non-linear phase'', ``Borel measure'', ``unequal
atoms'', ``spectral measure'', and ``exponentials''. They found the source
paper but no later explicit answer. This packet should therefore be reviewed
as a likely new, literal counterexample rather than as a literature retrieval.

\section{Scope and Limitations}

This answers the first open problem as literally stated for arbitrary finite
Borel measures. It does not address the remaining open problems in Section 8:
regular nonlinear phases for Cantor-type supports or open connected sets, the
one-dimensional continuous irregular phase problem for Lebesgue measure, or
classification of higher-dimensional $C^1$ phases on cubes.

\section{Human Review Recommendation}

Review should focus on whether the intended problem silently assumes nonatomic
or continuous measures. If not, the unequal two-atom measure gives a complete
negative answer.

\begin{thebibliography}{9}
\bibitem{GLO}
J.-P. Gabardo, C.-K. Lai, and V. Oussa,
\emph{On exponential bases and frames with non-linear phase functions and some
applications}, arXiv:2007.03736, 2020.
\end{thebibliography}

\end{document}
